\begin{table}[!htbp]\centering
\caption{Grupos educativos y exposición antes de la reforma}\label{tab:journal_desc}
\begin{threeparttable}\small
\setlength{\tabcolsep}{5pt}
\begin{tabular}{lrr}
\toprule
Indicador & Baja educación & Alta educación \\
\midrule
Observaciones: población de 14--65 años & 65954 & 27278 \\
Edad (años) & 35.9 & 34.6 \\
Zona urbana (\%) & 76.2 & 93.1 \\
Observaciones: muestra salarial & 14449 & 7795 \\
Mediana salarial mensual nominal (S/) & 946 & 1200 \\
Salario real por hora, media (S/) & 5.80 & 9.26 \\
Salario mensual inferior a S/930 (\%) & 45.2 & 25.2 \\
Salario mensual inferior a S/1025 (\%) & 56.9 & 40.0 \\
Salario mensual en [S/930, S/1025) (\%) & 11.7 & 14.9 \\
Observaciones: muestra de composición & 46090 & 15463 \\
Empleo formal (\%) & 11.6 & 38.3 \\
Cuenta propia (\%) & 42.8 & 30.9 \\
\bottomrule
\end{tabular}
\begin{tablenotes}\footnotesize
\item ENAHO 2021T1--2022T1, edades 14--65. Medias, proporciones y mediana ponderadas por fac500a; observaciones sin ponderar. Muestra salarial: asalariados privados con salario horario válido, incluidos trabajadores del hogar. Muestra de composición: ocupados no públicos. Los porcentajes bajo cada piso usan remuneración mensual equivalente nominal y no certifican incumplimiento legal: no se restringe a jornada completa. La franja entre pisos mide ubicación salarial previa, no asignación efectiva al tratamiento.
\end{tablenotes}\end{threeparttable}\end{table}
